% Sección A · Registro de mejora (máximo 1 página, va primero). Las referencias de sección y página
% se generan con \ref y \pageref: no escribirlas a mano.
\providecommand{\pend}[1]{\textcolor{red}{#1}}
\newcolumntype{L}{>{\raggedright\arraybackslash}X}
\newcolumntype{P}[1]{>{\raggedright\arraybackslash}p{#1}}

\section{Registro de mejora}
\label{sec:mejora}

Este informe es la Entrega~1 corregida, con tres secciones nuevas: arquitectura, prototipo y adopción. La tabla recoge cada observación de la Entrega~1 (E1) y de la Presentación~1 (P1), qué hicimos y dónde quedó el cambio.

\begin{center}
\footnotesize
\renewcommand{\arraystretch}{1.0}
\begin{tabularx}{\textwidth}{@{}r P{0.95cm} L L P{1.35cm}@{}}
\toprule
\# & \textbf{Origen} & \textbf{Observación} & \textbf{Qué hicimos} & \textbf{Dónde} \\
\midrule
1 & E1 &
Cajas y litros es un problema de gobierno de datos; la empresa no declaró el litro como su unidad. &
\textbf{Atendida.} El mandato enuncia la fractura: comercial habla de cajas, operaciones de litros (V). Medir por litro es decisión nuestra (I). El acta lleva una nota de corrección. &
§\ref{sec:mandato}, p.~\pageref{sec:mandato}; acta \\
\addlinespace[2pt]
2 & E1 y P1 &
La definición que mueve el número es la fecha de conversión a pesos, no los 9, 18 o 4,5~L por caja. &
\textbf{Atendida.} Moneda y fecha pasan al cuerpo (USD 11,2\,\%; CLP 10,1\,\% a factura; 6,3\,\% a cobranza) y se eligen en el prototipo. Litros por caja, al anexo~D. &
§\ref{sec:metrica}, p.~\pageref{sec:metrica} \\
\addlinespace[2pt]
3 & E1 y P1 &
Mostrar el 108\,\% sin las líneas corregidas (montos, monedas, botellas). &
\textbf{Atendida.} Con esas líneas sin corregir, el alza sería 17,6\,\% y la mezcla explicaría el 78,2\,\% (M): sigue siendo la causa. Se puede ver en el prototipo. &
§\ref{sec:metrica}, p.~\pageref{tab:sensibilidades} \\
\addlinespace[2pt]
4 & E1 y P1 &
CLP 2.961~MM valoriza el pasado; falta cuánto vale cada punto de gran reserva. &
\textbf{Atendida.} Cada punto de gran reserva en vez de varietal vale CLP 353~MM de margen bruto al año (338 con el dólar de la cobranza). Va en la primera lámina de la P2. &
§\ref{sec:negocio}, p.~\pageref{sec:negocio} \\
\addlinespace[2pt]
5 & E1 &
Las acciones no dicen cuánto cuestan. &
\textbf{Atendida.} Incentivo de CLP 5.000 por caja incremental (CLP 127,1~MM al año para $+2,0$\,pp) y reasignación de stock sin costo; dueño: Gerencia Comercial. &
§\ref{sec:negocio}, p.~\pageref{sec:negocio} \\
\addlinespace[2pt]
6 & P1 &
«Medir y fijar metas» no es una decisión. &
\textbf{Atendida.} Meta de $+2,0$\,pp de Gran Reserva en mercados con recomendación (CLP 706~MM/año de margen), decidida mensualmente por Gerencia Comercial. &
§\ref{sec:negocio}, p.~\pageref{sec:negocio} \\
\addlinespace[2pt]
7 & P1 &
No convertir divisas con un promedio del período. &
\textbf{Atendida.} Toda cifra en pesos se convierte línea por línea con el dólar de su fecha; ya no se usa el 946,34. &
§\ref{sec:negocio}, p.~\pageref{sec:negocio} \\
\addlinespace[2pt]
8 & P1 &
Falta claridad en la fecha de corte y en la unidad de tiempo. &
\textbf{Atendida.} Facturas 1-ene a 31-dic-2025, por semestre; declarado en cada figura, en el prototipo y en la bitácora. &
§\ref{sec:metrica}, p.~\pageref{sec:metrica} \\
\addlinespace[2pt]
9 & P1 &
Explicar el modelo relacional con más detención. &
\textbf{En el informe, sí;} en la P2, no: la única lámina técnica permitida es la arquitectura. &
§\ref{sec:modelo}, p.~\pageref{sec:modelo} \\
\addlinespace[2pt]
10 & E1 &
Profundizar la modularización hacia ETL o ELT. &
\textbf{Atendida.} Arquitectura por capas con sus tres decisiones y pruebas de calidad. &
§\ref{sec:arquitectura}, p.~\pageref{sec:arquitectura} \\
\bottomrule
\end{tabularx}
\end{center}

\vspace{-4pt}
\paragraph{El número, E1 frente a hoy.}

\begin{center}
\footnotesize\setlength{\tabcolsep}{5pt}
\renewcommand{\arraystretch}{1.05}
\begin{tabular}{@{}l r r l@{}}
\toprule
\textbf{Cifra} & \textbf{E1} & \textbf{E2} & \textbf{Cambio} \\
\midrule
Precio medio por litro, 1.\textsuperscript{er} sem. 2025 & 4,574 USD/L & 4,574 USD/L & Ninguno \\
Precio medio por litro, 2.\textsuperscript{o} sem. 2025 & 5,085 USD/L & 5,085 USD/L & Ninguno \\
Alza del precio medio & +11,2\,\% & +11,2\,\% & Ninguno \\
Parte del alza que es mezcla & 108\,\% & 108\,\% & Ninguno \\
Ventas fuera por calidad & 1,6\,\% & 1,6\,\% & Ninguno \\
Valor de la mezcla, CLP & 2.961 MM, venta 2.\textsuperscript{o} sem. & 353 MM/año, margen por punto & Ver nota \\
\bottomrule
\end{tabular}
\end{center}

\noindent La métrica no cambió: ninguna observación tocó su definición ni su universo, y la cadena reproduce las cifras de la E1. Cambia el valor en pesos: ahora es lo que vale cada punto de mezcla hacia adelante (fila~4), en margen y no en venta, y convertido línea por línea con el dólar de su fecha (fila~7).
